% GENERATED FILE. Edit canonical lessons, then run npm run generate:books.
% canonical-source-hash: fnv1a64:8428f4bb428edd6e
% canonical-lessons: UR-C346-khidmat, UR-C346-isar, UR-C346-lalach, UR-C346-bhalai, UR-C346-jatana, UR-R346-review-339-346

\chapter{Serving Others}
\label{ch:ur-serving-others}
\hlchaptermodality{346}

\begin{chapteropening}
\textbf{By the end of this chapter:} \emph{I can talk about selfless and selfish acts.}

Complete the last lesson of chapter 346: I can talk about selfless and selfish acts.
\end{chapteropening}

\section[\texorpdfstring{\ur{خدمت} (khidmat)}{خدمت (khidmat)}]{\ur{خدمت} (khidmat) — service to others}

\label{lesson:UR-C346-khidmat}



\begin{quote}
Before the new one: say the Urdu for a principle, then the Urdu for a creed.
\end{quote}

\subsection*{You'll want to know: \ur{خدمت}}
\textbf{\ur{خدمت}} — \emph{khidmat} — \textquotedblleft{}service to others\textquotedblright{}.

\begin{quote}
\ur{وہ} \ur{ہسپتال} \ur{میں} \ur{مریضوں} \ur{کی} \ur{خدمت} \ur{کرتی} \ur{ہے۔}
\end{quote}

\begin{quote}
\emph{voh haspatāl meṅ marīzoṅ kī khidmat kartī hai.}
\end{quote}

\begin{quote}
She looks after patients at the hospital.
\end{quote}

Feminine.

\begin{cousinweb}
\textbf{\ur{خدمت}} (\emph{khidmat}) is usually traced, through Persian, to Arabic \emph{khidma}, \textquotedblleft{}service\textquotedblright{}.
\end{cousinweb}

\begin{grammarlens}[title={What you've built}]
One more word for giving reasons and opinions.
\end{grammarlens}

\subsection*{Your turn}
\begin{itemize}
\raggedright
  \item \emph{Say it:} \emph{khidmat}
  \item \emph{Say it:} \emph{khidmat}, once more
  \item \emph{Say it:} \emph{voh haspatāl meṅ marīzoṅ kī khidmat kartī hai.} — She looks after patients at the hospital.
  \item \emph{Recall:} say the Urdu for a principle, then the Urdu for a creed, then say \emph{khidmat} again
\end{itemize}

\begin{tcolorbox}[breakable,colback=teal!4,colframe=teal!35!black,title={Before you move on}]
What does \ur{خدمت} mean? (\textquotedblleft{}Service to others\textquotedblright{}.) Say it once more.
\end{tcolorbox}
\section[\texorpdfstring{\ur{ایثار} (īsār)}{ایثار (īsār)}]{\ur{ایثار} (īsār) — selflessness; self-sacrifice}

\label{lesson:UR-C346-isar}



\begin{quote}
Before the new one: say the Urdu for a creed, then the Urdu for service to others.
\end{quote}

\subsection*{You'll want to know: \ur{ایثار}}
\textbf{\ur{ایثار}} — \emph{īsār} — \textquotedblleft{}selflessness; self-sacrifice\textquotedblright{}.

\begin{quote}
\ur{ماں} \ur{کا} \ur{ایثار} \ur{کوئی} \ur{نہیں} \ur{بھولتا۔}
\end{quote}

\begin{quote}
\emph{māṅ kā īsār koī nahīṅ bhūltā.}
\end{quote}

\begin{quote}
No one forgets a mother's selflessness.
\end{quote}

Masculine.

\begin{cousinweb}
\textbf{\ur{ایثار}} (\emph{īsār}) is usually traced, through Persian, to Arabic \emph{īthār}, \textquotedblleft{}putting another before oneself\textquotedblright{}.
\end{cousinweb}

\begin{grammarlens}[title={What you've built}]
One more word for giving reasons and opinions.
\end{grammarlens}

\subsection*{Your turn}
\begin{itemize}
\raggedright
  \item \emph{Say it:} \emph{īsār}
  \item \emph{Say it:} \emph{īsār}, once more
  \item \emph{Say it:} \emph{māṅ kā īsār koī nahīṅ bhūltā.} — No one forgets a mother's selflessness.
  \item \emph{Recall:} say the Urdu for a creed, then the Urdu for service to others, then say \emph{īsār} again
\end{itemize}

\begin{tcolorbox}[breakable,colback=teal!4,colframe=teal!35!black,title={Before you move on}]
What does \ur{ایثار} mean? (\textquotedblleft{}Selflessness; self-sacrifice\textquotedblright{}.) Say it once more.
\end{tcolorbox}
\section[\texorpdfstring{\ur{لالچ} (lālach)}{لالچ (lālach)}]{\ur{لالچ} (lālach) — greed}

\label{lesson:UR-C346-lalach}



\begin{quote}
Before the new one: say the Urdu for service to others, then the Urdu for selflessness.
\end{quote}

\subsection*{You'll want to know: \ur{لالچ}}
\textbf{\ur{لالچ}} — \emph{lālach} — \textquotedblleft{}greed\textquotedblright{}.

\begin{quote}
\ur{لالچ} \ur{بری} \ur{بات} \ur{ہے۔}
\end{quote}

\begin{quote}
\emph{lālach burī bāt hai.}
\end{quote}

\begin{quote}
Greed is a bad thing.
\end{quote}

Masculine.

\begin{cousinweb}
\textbf{\ur{لالچ}} (\emph{lālach}) is inherited: CDIAL lists the Hindi cousin \emph{lālac}, \textquotedblleft{}cupidity\textquotedblright{}, under Sanskrit \emph{lālitya}, \textquotedblleft{}amorousness\textquotedblright{} (CDIAL 11029); \textbf{\ur{لالچی}}, \textquotedblleft{}greedy\textquotedblright{}, is built on it.
\end{cousinweb}

\begin{grammarlens}[title={What you've built}]
One more word for giving reasons and opinions.
\end{grammarlens}

\subsection*{Your turn}
\begin{itemize}
\raggedright
  \item \emph{Say it:} \emph{lālach}
  \item \emph{Say it:} \emph{lālach}, once more
  \item \emph{Say it:} \emph{lālach burī bāt hai.} — Greed is a bad thing.
  \item \emph{Recall:} say the Urdu for service to others, then the Urdu for selflessness, then say \emph{lālach} again
\end{itemize}

\begin{tcolorbox}[breakable,colback=teal!4,colframe=teal!35!black,title={Before you move on}]
What does \ur{لالچ} mean? (\textquotedblleft{}Greed\textquotedblright{}.) Say it once more.
\end{tcolorbox}
\section[\texorpdfstring{\ur{بھلائی} (bhalāī)}{بھلائی (bhalāī)}]{\ur{بھلائی} (bhalāī) — welfare; the good (of others)}

\label{lesson:UR-C346-bhalai}



\begin{quote}
Before the new one: say the Urdu for selflessness, then the Urdu for greed.
\end{quote}

\subsection*{You'll want to know: \ur{بھلائی}}
\textbf{\ur{بھلائی}} — \emph{bhalāī} — \textquotedblleft{}welfare; the good (of others)\textquotedblright{}.

\begin{quote}
\ur{وہ} \ur{ہمیشہ} \ur{دوسروں} \ur{کی} \ur{بھلائی} \ur{سوچتا} \ur{ہے۔}
\end{quote}

\begin{quote}
\emph{voh hameshā dūsroṅ kī bhalāī sochtā hai.}
\end{quote}

\begin{quote}
He always thinks of the good of others.
\end{quote}

Feminine.

\begin{cousinweb}
\textbf{\ur{بھلائی}} (\emph{bhalāī}) is built on \emph{bhalā}, \textquotedblleft{}good\textquotedblright{}, which is inherited: CDIAL lists the Hindi cousin \emph{bhalā} under Sanskrit \emph{bhalla}, \textquotedblleft{}auspicious\textquotedblright{} (CDIAL 9408); \emph{-āī} makes the noun.
\end{cousinweb}

\begin{grammarlens}[title={What you've built}]
One more word for giving reasons and opinions.
\end{grammarlens}

\subsection*{Your turn}
\begin{itemize}
\raggedright
  \item \emph{Say it:} \emph{bhalāī}
  \item \emph{Say it:} \emph{bhalāī}, once more
  \item \emph{Say it:} \emph{voh hameshā dūsroṅ kī bhalāī sochtā hai.} — He always thinks of the good of others.
  \item \emph{Recall:} say the Urdu for selflessness, then the Urdu for greed, then say \emph{bhalāī} again
\end{itemize}

\begin{tcolorbox}[breakable,colback=teal!4,colframe=teal!35!black,title={Before you move on}]
What does \ur{بھلائی} mean? (\textquotedblleft{}Welfare; the good (of others)\textquotedblright{}.) Say it once more.
\end{tcolorbox}
\section[\texorpdfstring{\ur{جتانا} (jatānā)}{جتانا (jatānā)}]{\ur{جتانا} (jatānā) — to make known; to make a show of (a favour)}

\label{lesson:UR-C346-jatana}



\begin{quote}
Before the new one: say the Urdu for greed, then the Urdu for welfare.
\end{quote}

\subsection*{You'll want to know: \ur{جتانا}}
\textbf{\ur{جتانا}} — \emph{jatānā} — \textquotedblleft{}to make known; to make a show of (a favour)\textquotedblright{}.

\begin{quote}
\ur{اس} \ur{نے} \ur{اپنی} \ur{خدمت} \ur{کبھی} \ur{نہیں} \ur{جتائی۔}
\end{quote}

\begin{quote}
\emph{us ne apnī khidmat kabhī nahīṅ jatāī.}
\end{quote}

\begin{quote}
She never made a show of her service.
\end{quote}

Past, with \ur{نے} (he or she …): \ur{اس} \ur{نے} \ur{اپنا} \ur{حق} \ur{جتایا}, \emph{us ne apnā haq jatāyā}.

\begin{cousinweb}
\textbf{\ur{جتانا}} (\emph{jatānā}) is inherited: CDIAL lists the Hindi cousin \emph{jatānā}, \textquotedblleft{}to make known\textquotedblright{}, under Sanskrit \emph{jñaptá}, \textquotedblleft{}instructed\textquotedblright{} (CDIAL 5273).
\end{cousinweb}

\begin{grammarlens}[title={What you've built}]
One more word for giving reasons and opinions.
\end{grammarlens}

\subsection*{Your turn}
\begin{itemize}
\raggedright
  \item \emph{Say it:} \emph{jatānā}
  \item \emph{Say it:} \emph{jatānā}, once more
  \item \emph{Say it:} \emph{us ne apnī khidmat kabhī nahīṅ jatāī.} — She never made a show of her service.
  \item \emph{Recall:} say the Urdu for greed, then the Urdu for welfare, then say \emph{jatānā} again
\end{itemize}

\begin{tcolorbox}[breakable,colback=teal!4,colframe=teal!35!black,title={Before you move on}]
What does \ur{جتانا} mean? (\textquotedblleft{}To make known; to make a show of (a favour)\textquotedblright{}.) Say it once more.
\end{tcolorbox}
\section[(dialogue)]{Looking back: from Rights and Duties to Serving Others}

\label{lesson:UR-R346-review-339-346}



\begin{quote}
Say the Urdu for authority and justice.
\end{quote}

\subsection*{Your turn}
Say these aloud:

\begin{itemize}
\raggedright
  \item authority, justice, suitable, to act on, to fulfil
  \item a verdict, bail, an accused person, an arbiter, to testify
  \item a ruler, a policy, a proposal, the power to rule, to hold to be
  \item a vote, an election, democracy, a candidate, to impose
  \item a government department, an application, a post in an office, an instruction, to suggest
  \item equal, discrimination, favouritism, oppression, to provoke
  \item worth, faith, honour, a principle, a creed
  \item service to others, selflessness, greed, welfare, to make known
\end{itemize}

\begin{tcolorbox}[breakable,colback=teal!4,colframe=teal!35!black,title={Before you move on}]
Authority? (\textbf{\ur{اختیار}}, \emph{ikhtiyār}.) Justice? (\textbf{\ur{انصاف}}, \emph{insāf}.) Suitable? (\textbf{\ur{مناسب}}, \emph{munāsib}.) To act on? (\textbf{\ur{عمل} \ur{کرنا}}, \emph{amal karnā}.) To fulfil? (\textbf{\ur{ادا} \ur{کرنا}}, \emph{adā karnā}.) A verdict? (\textbf{\ur{فیصلہ}}, \emph{faislā}.) Bail? (\textbf{\ur{ضمانت}}, \emph{zamānat}.) An accused person? (\textbf{\ur{ملزم}}, \emph{mulzim}.) An arbiter? (\textbf{\ur{منصف}}, \emph{munsif}.) To testify? (\textbf{\ur{گواہی} \ur{دینا}}, \emph{gavāhī denā}.) A ruler? (\textbf{\ur{حکمران}}, \emph{hukmrān}.) A policy? (\textbf{\ur{پالیسی}}, \emph{pālīsī}.) A proposal? (\textbf{\ur{تجویز}}, \emph{tajvīz}.) The power to rule? (\textbf{\ur{اقتدار}}, \emph{iqtidār}.) To hold to be? (\textbf{\ur{ٹھہرانا}}, \emph{ṭhahrānā}.) A vote? (\textbf{\ur{ووٹ}}, \emph{voṭ}.) An election? (\textbf{\ur{انتخاب}}, \emph{intikhāb}.) Democracy? (\textbf{\ur{جمہوریت}}, \emph{jamhūriyat}.) A candidate? (\textbf{\ur{امیدوار}}, \emph{ummīdvār}.) To impose? (\textbf{\ur{تھوپنا}}, \emph{thopnā}.) A government department? (\textbf{\ur{محکمہ}}, \emph{mahkamā}.) An application? (\textbf{\ur{درخواست}}, \emph{darkhāst}.) A post in an office? (\textbf{\ur{عہدہ}}, \emph{uhdā}.) An instruction? (\textbf{\ur{ہدایت}}, \emph{hidāyat}.) To suggest? (\textbf{\ur{سجھانا}}, \emph{sujhānā}.) Equal? (\textbf{\ur{برابر}}, \emph{barābar}.) Discrimination? (\textbf{\ur{امتیاز}}, \emph{imtiyāz}.) Favouritism? (\textbf{\ur{جانبداری}}, \emph{jānibdārī}.) Oppression? (\textbf{\ur{ظلم}}, \emph{zulm}.) To provoke? (\textbf{\ur{اکسانا}}, \emph{uksānā}.) Worth? (\textbf{\ur{قدر}}, \emph{qadr}.) Faith? (\textbf{\ur{ایمان}}, \emph{īmān}.) Honour? (\textbf{\ur{عزت}}, \emph{izzat}.) A principle? (\textbf{\ur{اصول}}, \emph{usūl}.) A creed? (\textbf{\ur{عقیدہ}}, \emph{aqīdā}.) Service to others? (\textbf{\ur{خدمت}}, \emph{khidmat}.) Selflessness? (\textbf{\ur{ایثار}}, \emph{īsār}.) Greed? (\textbf{\ur{لالچ}}, \emph{lālach}.) Welfare? (\textbf{\ur{بھلائی}}, \emph{bhalāī}.) To make known? (\textbf{\ur{جتانا}}, \emph{jatānā}.)
\end{tcolorbox}
